%% results.tex -- submission version

\section{RQ1: Recovery-Path Source Provenance}
\label{sec:results}

RQ1 asks how recovery paths represent the origin and authority of failure
content. We report a qualitative multiple-case result, not a prevalence
estimate for the entire agent ecosystem. The source audit fetched all six
fixed revisions and matched every required code predicate.

\subsection{Cross-case Taxonomy}
\label{sec:rq1-taxonomy}

All six cases lack a source-origin label that survives into the recovery
representation. Three cases additionally preserve an error-status field
(AutoGen, OpenHands, and LangGraph), but an error status answers
\textquotedblleft{}did execution fail?\textquotedblright{} rather than
\textquotedblleft{}where did this content originate?\textquotedblright{}.
Four cases combine the missing source-origin label with a recovery-action
frame: Reflexion, MetaGPT, Aider, and AutoGen. OpenHands and LangGraph are
boundary cases because they expose error status but do not add the same
action-authorizing recovery frame.

\begin{table}[htbp]
\centering
\small
\caption{Source-level recovery-path taxonomy at fixed revisions. Source provenance means origin tracking, not an error-status flag.}
\label{tab:taxonomy}
\setlength{\tabcolsep}{2pt}
\begin{tabular}{p{1.8cm}p{1.5cm}p{1.5cm}p{1.5cm}p{2.0cm}p{4.0cm}}
\toprule
\textbf{Case} & \textbf{Prov.} & \textbf{Error} & \textbf{Action} & \textbf{Sig.} & \textbf{Fixed-revision evidence} \\
\midrule
Reflexion & No & No & Yes & Yes & Raw \texttt{feedback} placed in a unit-test-results recovery envelope. \\
MetaGPT & No & No & Yes & Yes & Test stderr placed under \textquotedblleft Console logs\textquotedblright{} in a role-playing debug prompt. \\
Aider & No & No & Yes & Yes & Edit failure text and target content followed by an instruction to reply with fixed versions. \\
AutoGen & No & Yes & Yes & Yes & \texttt{is\_error=True} is preserved internally, but the model-facing result has no source-origin label. \\
OpenHands & No & Yes & No & No & \texttt{ErrorObservation} and an error marker are converted to a user message without source-origin labeling. \\
LangGraph & No & Yes & No & No & Exception status is represented by \texttt{ToolMessage} with \texttt{status=error}; returned error strings can remain success messages. \\
\bottomrule
\end{tabular}
\end{table}

\paragraph{Reflexion.}
The retry path concatenates the raw test-runner output with the previous
implementation and reflection text:

\begin{lstlisting}[language=Python,caption={Reflexion retry path at commit \texttt{218cf0e}.},label={lst:wild-reflexion}]
message = (f"[previous impl]:\n{add_code_block(prev_func_impl)}\n\n"
           f"[unit test results from previous impl]:\n{feedback}\n\n"
           f"[reflection on previous impl]:\n{self_reflection}\n\n"
           f"[improved impl]:\n{func_sig}")
# feedback is raw test-runner output; no source-origin label is carried.
\end{lstlisting}

\paragraph{MetaGPT.}
At \texttt{v0.8.1}, \texttt{DebugError} places test output under a
role-playing prompt:

\begin{lstlisting}[language=Python,caption={MetaGPT debug prompt at \texttt{v0.8.1}.},label={lst:wild-metagpt}]
# Console logs
{logs}
---
Now you should start rewriting the code:
# logs is test-runner output and is not source-origin tagged.
\end{lstlisting}

\paragraph{Boundary cases.}
OpenHands stores \texttt{ErrorObservation} in state history and adds
\texttt{[Error occurred in processing last action]} when converting the
observation to a user message. LangGraph exposes an error status for
exceptions in \texttt{ToolMessage}, but a tool that returns an error-like
string can still produce a successful tool message. AutoGen's internal
\texttt{FunctionExecutionResult} correctly records \texttt{is\_error=True};
the limitation is that this status is not a source-origin distinction in the
model-facing recovery content. These observations are why the paper separates
source provenance from error status.

\subsection{Interpretation}
\label{sec:rq1-trampoline}

The four full-signature cases identify a recurring trust assumption: failure
content is treated as material for the next repair action without preserving
whether it originated in an external file, test fixture, tool return, or
framework-generated diagnostic. The source evidence supports a design-level
finding about missing complete mediation at the recovery boundary. It does not
support a statistical claim about all agent frameworks, and it does not by
itself establish exploitability; exploitability is evaluated in RQ2.
